\documentclass[11pt, a4paper]{article}
\usepackage[a4paper,
left=5mm,
right=5mm,
top=15mm,
bottom = 5mm]{geometry}
\usepackage{amsmath}
\usepackage{graphicx}
\usepackage{algorithm}
\usepackage{algorithmic}
\usepackage{multicol}
\usepackage{amsthm}
\usepackage{bm}
\usepackage{fancyhdr}

\newcommand\tab[1][1cm]{\hspace*{#1}}
\renewcommand{\labelitemii}{$\star$}

\begin{document}


\begin{titlepage}
	\begin{center}
		\vspace*{1cm}
		
		\Huge{Tesina\\Human Computer Interaction}
		\vspace{1.5cm}
		\Huge
		\textbf{\\CiakTime}
		\vspace{1.5cm}
		
		\Large
		Authors:\\
		\textbf{Mauro Ficorella 1941639}\\
		\textbf{Martina Turbessi 1944497}\\
		\textbf{Valentina Sisti 1952657}\\
		\vspace{0.5cm}
		
		\vfill
		
		\includegraphics[width=0.4\textwidth]{Images/Logo.jpg}
		
		\vfill
		
		\vspace{0.8cm}
		
		\Large
		Sapienza\\
		July 2021
	\end{center}
\end{titlepage}


% INTRODUZIONE --------------------------------------------------------------------

\newpage

\section*{Introduction}

The idea behind CiakTime comes from the fact that nowadays there are a lot of movies out 
both in theatres and in streaming platforms, so that cinema lovers can satisfy their
needs to stay updated with the latest movies and keep track of them. 
Also they may want to know on which streaming platforms they can found the movies.
Finally they could have the necessity to interact with other cinema lovers about their 
favourite movies or share their opinion about the movie with the community through reviews.
\\\\
For these reasons, our app offers a lot of functionalities. 
The user has the possibility to keep track of already watched movies, movies to watch and favourite movies;
moreover he can search for movies by title, also filtering results, search for actors and movie directors, 
look for upcoming movies and popular movies and actors.
Regarding the movies, he can read information about plot, cast, year of release, duration, genre,
movie director and streaming platform on which the movie is available; in addition, he can 
review and rate movies, comment and like reviews made by other users.
Finally, regarding movie directors and actors, the user can read t



% CAPITOLO 1 -------------------------------------------------------------------------

\newpage

\section*{NOME}

% CAPITOLO 2 -------------------------------------------------------------------------

\newpage

\section*{NOME}

% CAPITOLO 3 -------------------------------------------------------------------------

\newpage

\section*{NOME}

% CAPITOLO 4 -------------------------------------------------------------------------

\newpage

\section*{NOME}

% CAPITOLO 5 -------------------------------------------------------------------------

\newpage

\section*{NOME}

% CAPITOLO 6 -------------------------------------------------------------------------

\newpage

\section*{NOME}

% CAPITOLO 7 -------------------------------------------------------------------------

\newpage

\section*{NOME}

% CAPITOLO 8 -------------------------------------------------------------------------

\newpage

\section*{NOME}

\end{document}

% COSE UTILI --------------------------------------------------------------------------

%\section*{NOME}
%\subsection*{1.1}
%\setlength{\intextsep}{0pt} --> elimina lo spazio
%\vspace{-3mm}
%\hspace*{0cm}

% Font -------------------------------------

%GRASSETTO: \textbf

% Simboli ---------------------------------

%$\leftarrow$

% Elenco puntato ----------------------

%\begin{itemize}
%\setlength\itemsep{0.01em}
%\item 1
%\item 2
%\end{itemize}

% Graffa grande -----------------------

%\[  
%    \left\{ 
%    \begin{array}{ll} 
%      \mbox{1}
%      \mbox{2}
%    \end{array}
%    \mbox{riga al lato}
%   \right. 
%\]

% Multicolonne --------------------------

% \begin{multicols}{2}
% \columnbreak
% \end{multicols}

% Algoritmi -------------------------------

%\renewcommand{\thealgorithm}{1.\arabic{algorithm}}
%\setcounter{algorithm}{0}
%\begin{algorithm}
%\footnotesize
%\caption{Nome}
%\textbf{Input:} \\
%\textbf{Output:} 

%\begin{algorithmic}[1]
%\STATE 
%\FOR{ = 0 \TO i = n} ---- \ENDFOR
%\IF{} ---- \ELSIF{} ---- \ENDIF
%\RETURN 
%\end{algorithmic}
%\end{algorithm}

% Minipage ------------------------------

%\begin{minipage}[t]{0.5\textwidth}

% queste 3 righe vanno attaccate
%\end{minipage}
%\hspace{0.02\linewidth}
%\begin{minipage}[t]{0.47\textwidth} 

%\begin{minipage}[t]{0.3\textwidth} 
%\end{minipage}

% Proof --------------------------------
%\begin{proof}[\textbf{per cambiare nome}]
%\end{proof}

